% 2de graad, reeks Telproblemen.
% Klimt: de levensboom, een stam die zich telkens vertakt.
% Motief: een boomdiagram met drie keuzes uit twee, samen 2 x 2 x 2 = 8 mogelijkheden.
\def\omslagbestand{omslag-levensboom.jpg}
\def\omslagbron{Gustav Klimt, \textit{De levensboom}, ontwerp voor de Stocletfries, 1910-1911 (publiek domein)}
\def\omslagrand{10.6 + 0.5*sin(\x*360/21)}
\def\omslagmotief{%
  \randlijn
  \foreach \a in {-1,1}{%
    \draw[omslaglijn=1pt] (16.5,8.1) -- ({16.5+1.6*\a},9.8);
    \foreach \b in {-1,1}{%
      \draw[omslaglijn=1pt] ({16.5+1.6*\a},9.8) -- ({16.5+1.6*\a+.8*\b},11.5);
      \foreach \c in {-1,1}{%
        \draw[omslaglijn=1pt] ({16.5+1.6*\a+.8*\b},11.5) -- ({16.5+1.6*\a+.8*\b+.4*\c},13.2);
        \fill[fel,draw=accent,line width=.6pt] ({16.5+1.6*\a+.8*\b+.4*\c},13.2) circle (1.5mm);
      }
      \fill[omslagstip] ({16.5+1.6*\a+.8*\b},11.5) circle (1mm);
    }
    \fill[omslagstip] ({16.5+1.6*\a},9.8) circle (1mm);
  }
  \fill[omslagstip] (16.5,8.1) circle (1.3mm);
  \node[text=white,font=\small,inner sep=1pt,anchor=west] at (17,8.1) {$2\cdot 2\cdot 2=8$};
}
